\subsection{Use of AI/LLM tools}
\label{app:ai-tools}

Large language models were used throughout this project, following the course
guidelines\footnote{\url{https://github.com/EducationalMaterialUiO/MachineLearningUiO/blob/main/LLM_Usage_Declaration_Guidelines.md}},
which also define the contribution levels below. All code was tested, all
derivations were checked, and every claim in the report was read, understood
and verified by the author.

\subsubsection{Tools used}
\begin{itemize}
    \item Claude Code (Claude Opus 5.5, accessed via the command-line interface,
    November 2026)
    \item ChatGPT (various models, accessed via the web, November 2026)
\end{itemize}

Claude Code was used as a coding assistant with access to the repository: to build the
classification experiments, extend the training loop, write tests, and fix linting and type
errors. It was also used to review and restructure the report, and to give feedback on
drafted sections. ChatGPT was sometimes used to ask for feedback on sections of the
report and to discuss results.

\subsubsection{Text writing and editing}
\Cref{tab:llm-writing-contribution} gives the contribution level for each part
of the report.

\begin{table}[htbp]
    \centering
    \caption{LLM contribution to text writing and editing (levels 0--3).}
    \label{tab:llm-writing-contribution}
    {\footnotesize
    \begin{tabularx}{\columnwidth}{@{}lcX@{}}
        \toprule
        \textbf{Section} & \textbf{Level} & \textbf{Notes} \\
        \midrule
        Abstract & TODO & \\
        Introduction & TODO & Revised by LLM from the author's draft: made the opening
        specific to this project, stated the aim and datasets, and shortened the outline \\
        % Levels above 1 are not allowed by the course; set after rewriting the LLM drafts
        Methods & TODO & Section order restructured by LLM; drafts of the MSE, penalty,
        momentum, adaptive learning rate and classification metric parts, and a revision of
        the binary cross entropy part; shortened the part on output layers and loss functions;
        drafts of the weight initialization and regularization parts; corrected the
        backpropagation equations \\
        Implementation & TODO & Drafted and shortened by LLM \\
        Results and discussion & TODO & LLM suggested moving the MNIST description to
        \cref{sec:implementation}; drafted the regularization part of the regression results
        and revised two regression paragraphs; drafted the classification results \\
        Conclusion & TODO & \\
        Appendices & TODO & \\
        \bottomrule
    \end{tabularx}}
\end{table}

\subsubsection{Code generation and assistance}
\Cref{tab:code-generation} gives the level for each source file (levels 0--4).
Functions with level 2 or above carry an \texttt{LLM-assisted} tag in their
docstring.

\begin{table*}[htbp]
    \centering
    \caption{LLM contribution to the code (levels 0--4).}
    \label{tab:code-generation}
    {\footnotesize
    \begin{tabularx}{\textwidth}{@{}l c X@{}}
        \toprule
        \textbf{File} & \textbf{Level} & \textbf{Description} \\
        \midrule
        \texttt{regression.ipynb} & 3 & Refactoring code and markdown cells for readability
        and write the code for producing \cref{fig:model-comparison}. Parallelize the CV search for
        optimal hyperparameters. \\
        \texttt{classification.ipynb} & 4 & Wrote most of the notebook: the logistic regression
        baseline, the hyperparameter searches with cached results, the comparison of learning
        rate schedules, the figures, and the comparison with Scikit-Learn. Grids and experiment
        design chosen by the author. \\
        \texttt{nn/network.py} & 2 & Completed the author's draft of learning rate schedules
        and early stopping in \texttt{FFNN.fit}. \\
        \texttt{tests/test\_network.py} & 4 & Wrote the test comparing backpropagation with
        \texttt{jax.grad}. \\
        \texttt{tests/test\_logistic.py} & 1 & Fixed two type errors flagged by mypy. \\
        \bottomrule
    \end{tabularx}}
\end{table*}

\subsubsection{Declared and undeclared use}
TODO: one example of a use that must be declared, and one example of a use
that need not be declared (e.g. using an LLM as a tutor and then writing
about it in your own words).
